\documentclass[11pt]{article}
\usepackage[margin=1in]{geometry}
\usepackage{amsmath,amssymb}
\usepackage{booktabs}
\usepackage[dvipsnames]{xcolor}
\usepackage{hyperref}
\hypersetup{colorlinks=true,linkcolor=blue!50!black,urlcolor=blue!50!black,citecolor=blue!50!black}

\title{When Precision Misleads:\\
A Subgroup Audit of AlphaFold's pLDDT as a Predictor\\
of Curated Protein Disorder}
\author{David Lindgreen}
\date{October 9, 2026 \\[2pt] \small Version 0.1.0 \quad Preprint, not peer reviewed}

\emergencystretch=3em
\begin{document}
\maketitle

\begin{abstract}
AlphaFold's per-residue confidence score (pLDDT) is lower inside intrinsically disordered regions, and a pLDDT$<$70
rule is used as a disorder predictor. We tested this on 228{,}662 residues of 387 human proteins joined to curated
DisProt annotations, with hypotheses and thresholds fixed before analysis. The core result is robust: median pLDDT is
45.1 inside curated disorder and 88.8 outside, the direction holds in 332 of 372 proteins analysed individually, and
the rule has recall 0.765 but precision 0.341 (MCC 0.367). We then audit the subgroup comparisons in our own
original report, which scored each subgroup against the same shared pool of non-disorder residues. For such designs
precision and MCC depend mechanically on subgroup size: at equal recall, the conditional-folding group's precision
would be 0.071 against 0.306 for the rest, almost exactly the observed 0.063 against 0.310. Compared with
size-independent quantities and protein-clustered intervals, the classifier's recall is lower on proteins with a
conditional-folding annotation (0.681 against 0.780; difference 0.099, 95\% interval 0.0005 to 0.212), while the
AUC difference (0.049; $-0.005$ to 0.115) is not distinguishable from zero. A reported weakest evidence code
(hydrogen--deuterium exchange, recall 0.351) rests on four proteins (interval 0.09 to 0.72). The original claims are
corrected accordingly.
\end{abstract}

\section{Question}
Alderson et al.~\cite{alderson2023} showed that some conditionally folded disordered regions receive confident
AlphaFold2 predictions. We asked a narrower, independently checkable question~\cite{jumper2021}: on a fresh
snapshot of human DisProt entries, is pLDDT lower inside curated disorder, how well does the fixed rule pLDDT$<$70
recover the curated labels, and where does it fail? This is not a replication of the earlier paper's figures; we compare
only qualitatively.

\section{Data and preregistered analysis}
We sampled 400 of 1{,}339 human DisProt entries with a fixed seed, joined each to its AlphaFold DB per-residue pLDDT,
and applied preregistered exclusions (13 entries: 7 without a resolvable prediction, 6 with a fragmented one). This
leaves 387 proteins and 228{,}662 residues, of which 36{,}988 (16.2\%) lie in a curated disorder region and 191{,}674 do
not. Hypotheses: \textbf{H1}, pLDDT is lower inside curated disorder (Mann--Whitney, with Kolmogorov--Smirnov as a
robustness check); \textbf{H2}, the rule pLDDT$<$70 predicts the curated label above chance (precision, recall, F1,
Matthews correlation coefficient MCC~\cite{matthews1975}).

\section{Preregistered results}
\textbf{H1 is supported.} Median pLDDT is 45.09 inside and 88.75 outside curated disorder (difference 43.66, bootstrap
95\% interval 43.34--43.97; KS $D=0.484$). Because residues of one protein are correlated, we also compared
within-protein medians for the 372 proteins containing both kinds of residue: the median paired difference is 36.47
(protein-cluster bootstrap 95\% interval 34.29--40.07), positive in 332 proteins, negative in 39 and tied in one.
\textbf{H2 is confirmed above chance, with moderate performance.}
\begin{table}[h]
\centering\small
\begin{tabular}{lcc}
\toprule
 & Point estimate & Protein-cluster 95\% interval\\
\midrule
Precision & 0.341 & 0.294--0.390\\
Recall & 0.765 & 0.726--0.802\\
F1 & 0.472 & 0.422--0.521\\
MCC & 0.367 & 0.320--0.415\\
\bottomrule
\end{tabular}
\caption{The rule pLDDT$<$70 against curated disorder. Confusion counts: TP 28{,}306, FP 54{,}760, TN 136{,}914, FN 8{,}682.}
\end{table}

\section{The flaw in the original subgroup comparisons}
The original report broke performance down by a protein-level conditional-folding flag (61 of 387 proteins carry a
``disorder to order'' annotation) and by evidence code, and called precision and MCC ``substantially worse'' for the
conditional-folding group (precision 0.063 against 0.310; MCC 0.142 against 0.360). Each row, however, scored its
subgroup's disorder residues against the \emph{same} 191{,}674 non-disorder residues, so the number of false positives
(54{,}760) is fixed while the number of true positives scales with subgroup size. Precision
$=\mathrm{TP}/(\mathrm{TP}+\mathrm{FP})$ therefore falls for any smaller subgroup, whatever the classifier does.
If both groups had the overall recall of 0.765, precision would be 0.071 for the 5{,}435 conditional-group residues and
0.306 for the other 31{,}553: the observed 0.063 and 0.310 are almost exactly this size effect. Precision and MCC are
simply not comparable across such subgroups. The original report noted this for evidence codes but not for the
conditional-folding table.

\section{Post-hoc size-independent audit}
We compared the groups using (i) recall within the group, and (ii) the area under the ROC curve (AUC) separating the
group's disorder residues from the shared non-disorder pool, with 95\% intervals from a protein-cluster bootstrap
(10{,}000 resamples for recall, 1{,}000 for AUC)~\cite{field2007}. Each protein belongs to exactly one group (61
conditional proteins against 326 others). This analysis was not preregistered and all of it is reported.
\begin{table}[h]
\centering\small
\begin{tabular}{lccc}
\toprule
 & Conditional & Other & Difference (95\% interval)\\
\midrule
Recall & 0.681 & 0.780 & 0.099 (0.0005 to 0.212)\\
AUC & 0.730 & 0.780 & 0.049 ($-$0.005 to 0.115)\\
Observed precision & 0.063 & 0.310 & (size-dependent)\\
Precision at equal recall & 0.071 & 0.306 & (size-dependent)\\
\bottomrule
\end{tabular}
\caption{Conditional-folding subgroup audit. Only recall and AUC are comparable across groups.}
\end{table}
The direction agrees with the motivation of Alderson et al.: the rule is somewhat less sensitive where conditional
folding occurs. But the recall interval barely excludes zero and the AUC interval includes it, so the evidence for a
gap is weak at this sample size, not ``substantial''.

\paragraph{Evidence codes.}
The original report singled out hydrogen--deuterium exchange mass spectrometry (HDX-MS) as the weakest evidence code
(recall 0.351). Those 578 residues come from only four proteins, and the protein-clustered interval is 0.09 to 0.72.
Several other codes rest on one to five proteins (for example dynamic light scattering: 384 residues from one protein), for
which bootstrap intervals are degenerate or unreliable. Recall ranges across codes with larger protein counts are
narrower and more informative: X-ray missing-residue evidence 0.702 (202 proteins; 0.622--0.781), NMR 0.772
(108 proteins; 0.708--0.828), circular dichroism 0.757 (37 proteins; 0.649--0.845).

\section{Limitations}
One DisProt snapshot, human proteins only, a 400-of-1{,}339 sample, and one threshold. The conditional-folding flag is a
protein-level proxy, not the region-level definition of the earlier work. Absence of a curated disorder annotation does
not mean order, so false positives are partly missing curation. Subgroups are overlapping annotations, so no causal
role of any experimental technique is implied. The audit is post-hoc and descriptive; it corrects over-reading in the
original report and does not test a new biological hypothesis.

\section*{Reproducibility and AI assistance}
Code, frozen registry, tests and the subgroup audit:\\
\url{https://github.com/lindgreendavid/foldings-edge}\\
Interactive laboratory:\\
\url{https://foldings-edge-interactive.lindgreendavid.workers.dev}\\
The frozen registry regenerates identically from the committed residue table. The analysis code, tests and drafts of this
text were produced with substantial assistance from an AI system (Claude, Anthropic); the author is responsible for the
content.

\begin{thebibliography}{9}
\bibitem{alderson2023} T.~R. Alderson, I.~Priti\v{s}anac, {\DJ}.~Kolari\'c, A.~M. Moses, J.~D. Forman-Kay, Systematic
identification of conditionally folded intrinsically disordered regions by AlphaFold2, \emph{Proc. Natl. Acad. Sci. USA}
\textbf{120}(44), e2304302120 (2023).
\bibitem{jumper2021} J.~Jumper et al., Highly accurate protein structure prediction with AlphaFold,
\emph{Nature} \textbf{596}, 583--589 (2021).
\bibitem{matthews1975} B.~W. Matthews, Comparison of the predicted and observed secondary structure of T4 phage lysozyme,
\emph{Biochim. Biophys. Acta} \textbf{405}, 442--451 (1975).
\bibitem{field2007} C.~A. Field and A.~H. Welsh, Bootstrapping clustered data,
\emph{J. R. Stat. Soc. B} \textbf{69}, 369--390 (2007).
\end{thebibliography}
\end{document}
